\documentclass[12pt,aspectratio=169]{beamer}
\input{header_beamer}
\usetheme{Boadilla}
\usecolortheme{beaver}
\setbeamertemplate{page number in head/foot}{}
\setbeamertemplate{navigation symbols}{}

\title{Lecture-04: Random Variables}
\date{}

\begin{document}
\pagestyle{empty}
\maketitle

\begin{frame}{Random Variable}
\begin{Definition}[Random variable]<1->
Consider a probability space $(\Omega, \sF, P)$. 
A \textbf{random variable} $X: \Omega \to \R$ is a real-valued function from the sample space to real numbers, 
such that for each $x \in \R$ %the event space $\sF$ is generated by 
the event 
\EQ{
A_X(x) \triangleq \set{\omega \in \Omega: X(\omega) \le x}
= \set{X \le x} = X^{-1}(-\infty, x] = X^{-1}(B_x) \in \sF.
}  
We say that the random variable $X$ is $\sF$-measurable.  
\end{Definition}
\begin{block}{Reminder}<2->
Recall that the set $A_X(x)$ is always a subset of sample space $\Omega$ for any mapping $X:\Omega\to\R$, 
and $A_X(x)\in\sF$ is an event when $X$ is a random variable.
\end{block}
\end{frame}
\begin{frame}{Random Variable}
\begin{Example}[Constant function] 
Consider a mapping $X:\Omega\to \set{c} \subseteq \R$ defined on an arbitrary probability space $(\Omega, \sF, P)$, 
such that $X(\omega) = c$ for all outcomes $\omega \in \Omega$. 
We observe that 
\EQ{
A_X(x) = X^{-1}(B_x) = 
\begin{cases}
\emptyset, & x < c,\\
\Omega, & x \ge c.
\end{cases}
}
That is $A_X(x) \in \sF$ for all event spaces, and hence $X$ is a random variable and measurable for all event spaces. 
\end{Example}
\end{frame}
\begin{frame}{Random Variable}
\begin{Example}[Indicator function] 
For an arbitrary probability space $(\Omega,\sF, P)$ and an event $A \in \sF$, 
consider the indicator function $\Ind{A}:\Omega \to [0,1]$. 
Let $x \in \R$, and $B_x = (-\infty, x]$, then it follows that 
\EQ{
A_X(x) = \Ind{A}^{-1}(B_x) 
= \begin{cases}
\Omega, & x \ge 1,\\
A^c, & x \in [0, 1),\\
\emptyset, & x < 0.
\end{cases} 
}  
That is, $A_X(x) \in \sF$ for all $x \in \R$, and hence the indicator function $\Ind{A}$ is a random variable.  
\end{Example}
\end{frame}

\begin{frame}{Random Variable}
\begin{block}{Reminder} 
\begin{itemize}
\item <1-> Since any outcome $\omega \in \Omega$ is random, so is the real value $X(\omega)$. % \in \R$. 
\item <2-> Probability is defined only for events and not for random variables. 
The events of interest for random variables are the \textbf{lower level sets} $A_X(x) = \set{\omega: X(\omega) \le x} = X^{-1}(B_x)$ for any real $x$. % \in \R$. 
\item <3-> Consider a probability space $(\Omega,\sF,P)$ and a random variable $X:\Omega \to \R$ that is $\sG$ measurable for $\sG \subseteq \sF$. 
If $\sG \subseteq \sH$, then $X$ is also $\sH$ measurable.
\end{itemize}
\end{block}
\end{frame}

\begin{frame}{Distribution function for a random variable}
\begin{Definition}[Distribution function]<1->
For an $\sF$ measurable random variable $X:\Omega\to\R$ defined on the probability space $(\Omega,\sF,P)$, 
we can associate a \textbf{distribution function} (CDF) $F_X: \R \to [0,1]$ 
such that for all $x \in \R$, 
\EQ{
F_X(x) \triangleq P(A_X(x)) = P(\set{X \le x}) = P \circ X^{-1}(-\infty, x] = P\circ X^{-1}(B_x).
}
\end{Definition}


\begin{Example}[Constant random variable]<2->
\begin{itemize}
\item Let $X:\Omega \to \set{c} \subseteq \R$ be a constant random variable defined on $(\Omega, \sF, P)$. \\
\item $F_X(x) = \Ind{[c,\infty)}(x)$ 
\item $P(\set{X = c}) = 1$. 
\end{itemize}
\end{Example}
\end{frame}
\begin{frame}{Distribution function for a random variable}
\begin{Example}[Indicator random variable] 
For an indicator random variable $\Ind{A}:\Omega \to \set{0,1}$ defined on a probability space $(\Omega, \sF, P)$ and an event $A \in \sF$,  
we have 
\EQ{
F_X(x) = 
\begin{cases}
1, & x \ge 1,\\
1- P(A), & x \in [0, 1),\\
0, & x < 0.
\end{cases}
}
\end{Example}
\end{frame}

\begin{frame}{Properties of distribution function for a random variable}
\begin{block}{Lemma : Properties of distribution function} 
The distribution function $F_X$ for any random variable $X$ satisfies the following properties. 
\begin{enumerate}
\item <1-> The distribution function is monotonically non-decreasing in $x \in \R$. 
\item <2-> The distribution function is right-continuous at all points $x \in \R$.
\item <3-> The upper limit is $\lim_{x \to \infty}F_X(x) = 1$ and the lower limit is $\lim_{x \to -\infty}F_X(x) = 0$. 
\end{enumerate}
\end{block}
\end{frame}
\begin{frame}{Properties of distribution function}
\begin{Proof}
Let $X$ be a random variable defined on $(\Omega, \sF, P)$. 
\begin{itemize}
\item <1-> Let $x_1, x_2 \in \R$ such that $x_1 \le x_2$.\\  
Then for any $\omega \in A_{x_1}$, $X(\omega) \le x_1 \le x_2$, and $\omega \in A_{x_2}$ \implies $A_{x_1} \subseteq A_{x_2}$. 
 
\item <2-> For any $x \in \R$, consider any monotonically decreasing sequence $x\in\R^\N$ such that $\lim_n x_n = x_0$.\\
$\big(A_{x_n} = X^{-1}(-\infty, x_n] \in \sF: n \in \N\big)$, is monotonically decreasing and hence $\lim_{n \in \N}A_{x_n} = \cap_{n \in \N}A_{x_n} = A_{x_0}$.  
\EQ{
F_X(x_0) = P(A_{x_0}) = P(\lim_{n \in \N}A_{x_n}) = \lim_{n \in \N}P(A_{x_n}) = \lim_{x_n \downarrow x}F(x_n). 
}
\end{itemize}
\end{Proof}
\end{frame}
\begin{frame}{Properties of distribution function}
\begin{block}{Proof Continued}
\begin{itemize}
\item <1-> Consider a monotonically increasing sequence $x\in\R^\N$ such that $\lim_{n}x_n = \infty$\\
Then $(A_{x_n} \in \sF: n \in \N)$ is a monotonically increasing sequence of sets and $\lim_nA_{x_n} = \cup_{n \in \N}A_{x_n} = \Omega$.\\
\EQ{
\lim_{x_n \to \infty}F_X(x_n) = \lim_nP(A_{x_n}) = P(\lim_n A_{x_n}) = P(\Omega) = 1.
}
\item <2-> Similarly, take a monotonically decreasing sequence $x\in \R^\N$ such that $\lim_{n}x_n = -\infty$,\\
Then $(A_{x_n} \in \sF: n \in \N)$ is a monotonically decreasing sequence of sets and $\lim_nA_{x_n} = \cap_{n \in \N}A_{x_n} = \emptyset$.\\
From the continuity of probability, $\lim_{x_n \to -\infty}F_X(x_n) = 0$. 
\end{itemize}
\end{block}
\end{frame}
\begin{frame}{Properties of distribution function}
\begin{block}{Reminder}
If two reals $x_1 < x_2$ then $F_X(x_1) \le F_X(x_2)$ with equality if and only if $P\set{(x_1 < X \le x_2}) = 0$. 
This follows from the fact that $A_{x_2} = A_{x_1}\cup X^{-1}(x_1, x_2]$. 
\end{block}
\end{frame}
\begin{frame}{Event space generated by a random variable}
\begin{Definition}[Event space generated by a random variable] 
\begin{itemize}
\item <1-> Let $X: \Omega \to \R$ be an $\sF$ measurable random variable defined on the probability space $(\Omega, \sF, P)$.
\item <2-> The smallest event space generated by the events $A_X(x) = X^{-1}(B_x) = X^{-1}(-\infty, x]$ for $x \in \R$ is called the \textbf{event space generated} by this random variable $X$, 
\item <3-> $\sigma(X) \triangleq \sigma(\set{A_X(x): x \in \R})$. 
\end{itemize}
\end{Definition}
\end{frame}
\begin{frame}{Event space generated by a random variable}
\begin{block}{Reminder}
\begin{itemize}
\item <1-> $\sigma(X) = \set{X^{-1}(B): B \in \sB(\R)}$.
\item <2-> $A_X(x) = X^{-1}(B_x)$  
\item <3-> If $B \in \sB(\R) = \sigma(\set{B_x:x\in \R})$, then $X^{-1}(B) \in \sigma(\set{A_X(x): x \in \R})$. 
\item <4-> If $A \in \sigma(X) = \sigma(\set{A_X(x): x \in \R})$, then $A = X^{-1}(B)$ for some $B \in \sigma(\set{B_x: x \in \R})$.  
\end{itemize}
\end{block}
\end{frame}
\begin{frame}{Event space generated by a random variable}
\begin{Example}[Constant random variable]<1->
Let $X:\Omega \to \set{c} \subseteq \R$ be a constant random variable $(\Omega,\sF,P)$. \\
The smallest event space generated by this random variable is 
\EQ{\sigma(X) = \set{\emptyset, \Omega}}. 
\end{Example}
\begin{Example}[Indicator random variable]<2->
Let $\Ind{A}$ be an indicator random variable defined $(\Omega,\sF,P)$ and event $A \in \sF$\\ 
The smallest event space generated by this random variable is 
\EQ{\sigma(X) = \sigma(\set{\emptyset, A^c, \Omega}) = \set{\emptyset, A, A^c, \Omega}} 
\end{Example}
\end{frame}

\begin{frame}{Discrete random variables}
\begin{Definition}[Discrete random variables]
\begin{itemize}
\item <1-> If $X: \Omega \to \sX \subseteq \R$ takes countable values on real-line, then it is called a \textbf{discrete random variable}. Range$\sX$ is countable.
\item <2->  $X$ is completely specified by the \textbf{probability mass function} 
\EQ{
P_X(x) = P(\set{X = x}), \text{ for all }x \in \sX. 
}
\end{itemize}
\end{Definition}
\end{frame}
\begin{frame}{Discrete random variables}
\begin{Example}[Bernoulli random variable] 
\begin{itemize}
    \item <1-> For $(\Omega, \sF, P)$, the \textbf{Bernoulli random variable} is $X: \Omega \to \set{0,1}$ and $P_X(1) = p$.
    \item <2-> $A \triangleq X^{-1}\set{1}$, and $P(A) = p$. 
    \item <3-> The distribution function $F_X$ is given by 
\EQ{
F_X = (1-p)\Ind{[0, 1)} + \Ind{[1, \infty)}.
}
\end{itemize}
\end{Example}
\end{frame}

\begin{frame}{Discrete random variables}
\begin{block}{Lemma}<1->
Any discrete random variable is a linear combination of indicator function over a partition of the sample space.
\end{block}

\begin{Proof}<2->
\begin{itemize}
    \item <2-> Define events $E_x \triangleq \set{\omega \in \Omega: X(\omega) = x} \in \sF$ for each $x \in \sX$. 
    \item <3-> The mutually disjoint sequence of events $(E_x \in \sF: x \in \sX)$ partitions the sample space $\Omega$.
            \EQ{
            X(\omega) = \sum_{x \in \sX}x\Ind{E_x}(\omega).
            }
\end{itemize}
\end{Proof}
\end{frame}
\begin{frame}{Simple random variables}
\begin{Definition} 
Any discrete random variable $X:\Omega \to \sX \subseteq\R$ defined over a probability space $(\Omega,\sF,P)$, 
with finite range is called a simple random variable. 
\end{Definition}
\end{frame}

\begin{frame}{Simple random variables}
\begin{Example}[Simple random variables]
\begin{itemize}
\item <1-> Let $X = \sum_{x \in \sX}x\Ind{A_X(x)}$ be a simple random variable on $(\Omega,\sF,P)$
\item <2-> $(A_X(x) = X^{-1}\set{x} \in \sF: x \in \sX)$ is a finite partition of $\Omega$.\\
\item <3-> $\sX  = \set{x_1, \dots, x_n}$ where $x_1 \le \dots \le x_n$. 
\item <4->
\EQ{
X^{-1}(-\infty, x] = \begin{cases}
\Omega, &x \ge x_n,\\
\cup_{j=1}^iA_X(x_j), & x \in [x_i, x_{i+1}), i \in [n-1],\\
\emptyset, &x < x_1.
\end{cases}
}
\item <5-> The smallest event space generated by $X$ is $\set{\cup_{x \in S}A_X(x): S \subseteq \sX}$. 
\end{itemize}
\end{Example}
\end{frame}

\begin{frame}{Continuous random variables}
\begin{Definition}[Density function]<1->
For a continuous random variable $X$, 
there exists \textbf{density function} $f_X: \R \to [0, \infty)$ such that 
\EQ{
F_X(x) = \int_{-\infty}^xf_X(u)du.
}
\end{Definition}


\begin{Example}[Gaussian random variable]<2->
For a probability space $(\Omega, \sF, P)$, \textbf{Gaussian random variable} is a continuous random variable $X: \Omega \to \R$ defined by its density function 
\EQ{
f_X(x) = \frac{1}{\sqrt{2\pi}\sigma}\exp\left(-\frac{(x-\mu)^2}{2\sigma^2}\right),~ x \in \R.
}
\end{Example}
\end{frame}
\end{document}
